% ===== PRÉAMBULE CANONIQUE — à coller VERBATIM en tête de CHAQUE .tex =====
\documentclass[11pt,a4paper]{article}
\usepackage[utf8]{inputenc}\usepackage[T1]{fontenc}\usepackage{lmodern}
\usepackage[french]{babel}
\usepackage{amsmath,amssymb,mathtools}\usepackage{icomma}
\usepackage{siunitx}\sisetup{locale=FR}  % \qty{}{}, \si{}, \ang{}
\usepackage{xcolor}\usepackage{colortbl}
\usepackage{tikz}\usepackage{pgfplots}\pgfplotsset{compat=1.18}
\usepackage[siunitx,europeanresistors]{circuitikz} % circuits (élec.)
\usepackage[most]{tcolorbox}\usepackage{enumitem}\usepackage{array,booktabs}
\usepackage[a4paper,margin=2.2cm]{geometry}
\usetikzlibrary{arrows.meta,calc,angles,quotes,positioning,patterns,%
  backgrounds,shapes.geometric,decorations.pathmorphing,decorations.markings,intersections}
\definecolor{primary}{RGB}{222,49,99}\definecolor{primarydark}{RGB}{150,22,55}
\definecolor{defcol}{RGB}{0,114,178}\definecolor{methcol}{RGB}{52,92,148}
\definecolor{excol}{RGB}{0,135,90}\definecolor{attncol}{RGB}{200,40,55}
\tcbset{boxbase/.style={enhanced,breakable,boxrule=0.9pt,arc=2.5pt,
  left=3mm,right=3mm,top=2.3mm,bottom=2.3mm,fonttitle=\bfseries,coltitle=white}}
% --- boîtes TUTORIEL (noms EXACTS, ne pas renommer) ---
\newtcolorbox{cle}[1][]{boxbase,colback=defcol!6,colframe=defcol,
  title={\textsc{À retenir}\ifx\relax#1\relax\relax\else\textmd{ : \itshape #1}\fi}}
\newtcolorbox{methode}[1][]{boxbase,colback=methcol!7,colframe=methcol,sharp corners,
  title={\textsc{Méthode}\ifx\relax#1\relax\relax\else\textmd{ --- \itshape #1}\fi}}
\newtcolorbox{exemplebox}[1][]{boxbase,colback=excol!5,colframe=excol,
  title={\textsc{Exemple}\ifx\relax#1\relax\relax\else\textmd{ : \itshape #1}\fi}}
\newtcolorbox{piege}[1][]{boxbase,colback=attncol!6,colframe=attncol,
  title={\textsc{Piège}\ifx\relax#1\relax\relax\else\textmd{ : \itshape #1}\fi}}
\newtcolorbox{remarque}[1][]{boxbase,colback=black!4,colframe=black!45,
  title={\textsc{Remarque}\ifx\relax#1\relax\relax\else\textmd{ : \itshape #1}\fi}}
% --- boîtes EXERCICE / CORRIGÉ (noms EXACTS, auto-comptées) ---
\newtcolorbox[auto counter]{exo}[1][]{boxbase,colback=primary!4,colframe=primary,
  title={\textsc{Exercice}~\thetcbcounter\ifx\relax#1\relax\relax\else\textmd{ --- \itshape #1}\fi}}
\newtcolorbox[auto counter]{corrige}[1][]{boxbase,colback=excol!4,colframe=excol,
  title={\textsc{Corrigé}~\thetcbcounter\ifx\relax#1\relax\relax\else\textmd{ --- \itshape #1}\fi}}
\newcommand{\vv}[1]{\vec{#1}}\newcommand{\ds}{\displaystyle}
\newcommand{\R}{\mathbb{R}}\newcommand{\N}{\mathbb{N}}\newcommand{\Z}{\mathbb{Z}}
% =========================================================================
\begin{document}
\usetikzlibrary{babel}
\begin{exo}[à quelle température la résistance double-t-elle ?]
Un fil d'aluminium a une résistance $R_0=\qty{50}{\ohm}$ à $\qty{0}{\degreeCelsius}$ et un coefficient
de température $\alpha=\qty{4e-3}{\per\degreeCelsius}$.
\begin{enumerate}[label=\alph*)]
  \item À quelle température $T$ sa résistance est-elle \emph{doublée} ($R=\qty{100}{\ohm}$) ?
  \item À quelle température serait-elle \emph{triplée} ?
  \item Ces valeurs sont-elles à prendre au pied de la lettre ? (penser à la validité du modèle)
\end{enumerate}
\end{exo}

\begin{exo}[trois composés : trouver l'affirmation fausse]
Trois composés ont pour résistivités $\rho_1=\qty{1.6e-8}{\ohm\metre}$,
$\rho_2=\qty{22e-8}{\ohm\metre}$ et $\rho_3=\qty{3e6}{\ohm\metre}$. On en fait des fils de même
longueur et de même section. Indique si chaque affirmation est \textbf{vraie} ou \textbf{fausse} :
\begin{enumerate}[label=\alph*)]
  \item le composé 3 peut être un isolant ;
  \item parmi les trois, le composé 1 est le meilleur conducteur électrique ;
  \item la conductivité du composé 1 est supérieure à celle du composé 2 ;
  \item la conductivité du composé 3 est supérieure à celle du composé 1.
\end{enumerate}
\end{exo}

\begin{exo}[dimensionner une section]
On veut qu'un fil de longueur $L=\qty{10}{\metre}$ et de résistivité $\rho=\qty{22e-8}{\ohm\metre}$ ait
une résistance $R=\qty{2}{\ohm}$.
\begin{enumerate}[label=\alph*)]
  \item Quelle section $S$ faut-il lui donner ? Exprime-la en \si{\metre\squared} puis en
        \si{\milli\metre\squared}.
  \item Vrai ou faux : « la section doit valoir $\qty{2.2}{\milli\metre\squared}$ ». Justifie.
\end{enumerate}
\end{exo}

\begin{exo}[câble de cuivre : synthèse]
Un câble de cuivre ($\rho=\qty{1.7e-8}{\ohm\metre}$) a une longueur $L=\qty{100}{\metre}$ et une
section $S=\qty{2}{\milli\metre\squared}$.
\begin{enumerate}[label=\alph*)]
  \item Calcule sa résistance $R$ à $\qty{0}{\degreeCelsius}$.
  \item Pour diviser cette résistance par 2 sans changer ni le matériau ni la longueur, quelle
        nouvelle section faut-il ?
  \item En service, le câble s'échauffe à $\qty{50}{\degreeCelsius}$ ($\alpha=\qty{3.9e-3}{\per\degreeCelsius}$).
        Que devient la résistance de la question a) ?
\end{enumerate}
\end{exo}

\end{document}
